\begin{abstract}
Reinforcement-learning traders often settle on supracompetitive outcomes
without communicating: in Kyle markets, Q-learning informed traders trade far
less than competition implies, earning near-cartel profits and making prices
less informative. Whether this is \emph{collusion}, an outcome sustained by
the threat of punishment, or a \emph{learning bias} that prunes aggressive
actions matters for policy: the first calls for limits on transparency, the
second for standards on algorithm design. We develop diagnostics that separate
the two: a paired rival-deviation test, a noise-shock test, a no-memory
control, learning-rate sensitivity, and a new \emph{myopic placebo} that
retrains the same learners with no concern for the future and therefore no
reason to punish. We show analytically that in the standard Kyle market a
one-period deviation moves order flow by only $(I-1)/(2I)$ noise standard
deviations per unit of the value signal, whatever the volume of noise
trading, while with a large mass of information-insensitive investors the
ratio is in the hundreds, which explains why punishment-based collusion has
been reported only in that regime. In simulations with hundreds of
independent markets per configuration, Q-learners in the standard Kyle market
reach up to \ResidDeltaInt{} of the way from competition to the cartel, yet
rivals never punish deviations, even with perfect monitoring of each other's
orders, and deviating always pays. \AbstractPlaceboSentence{} We release a
validated, CPU-only testbed with exact benchmarks, tabular and deep RL agents,
and all diagnostics.
\end{abstract}
